% ---------------------------------------------------------------------------
% Proyecto I - Robot Aspiradora Autonomo (Yocto)
% CE-1113 Sistemas Empotrados - TEC
%
% Documento de Diseno (DI): indicadores DI1 a DI4 del enunciado.
% El detalle tecnico (calculos, listas de materiales, procedimientos) vive en
% docs/; aqui se resume y se cita.
%
% Compilar:  pdflatex documento-diseno.tex   (dos veces)
% ---------------------------------------------------------------------------
\documentclass[11pt,a4paper]{article}

\usepackage[T1]{fontenc}
\usepackage[utf8]{inputenc}
\usepackage{lmodern}
\usepackage[spanish,es-noshorthands,es-noquoting]{babel}
\usepackage[margin=2.4cm]{geometry}
\usepackage{tikz}
\usepackage{amsmath}
\usepackage{amssymb}
\usepackage{booktabs}
\usepackage{tabularx}
\usepackage{longtable}
\usepackage{array}
\usepackage{enumitem}
\usepackage{caption}
\PassOptionsToPackage{hyphens}{url}
\usepackage[hidelinks]{hyperref}

\usetikzlibrary{positioning,arrows.meta,fit,backgrounds,calc}

\captionsetup{font=small,labelfont=bf,width=.92\linewidth}
\setlength{\parskip}{0.45em}
\setlength{\parindent}{0pt}
\emergencystretch=4em
\setlist{itemsep=1pt,topsep=2pt}
\renewcommand{\arraystretch}{1.18}

% --- paleta (la misma de hardware-diagramas.tex) ----------------------------
\definecolor{cLog}{RGB}{31,102,168}
\definecolor{cPot}{RGB}{193,84,26}
\definecolor{cGris}{RGB}{120,120,120}
\definecolor{cFondo}{RGB}{245,247,250}
\definecolor{cPend}{RGB}{176,32,32}

\tikzset{
  blk/.style   ={draw,rounded corners=2pt,align=center,inner sep=4pt,
                 minimum height=9mm,font=\small},
  blkl/.style  ={blk,draw=cLog,fill=cLog!8},
  blkp/.style  ={blk,draw=cPot,fill=cPot!8},
  blkn/.style  ={blk,draw=cGris,fill=black!3},
  ln/.style    ={-{Latex[length=2mm]},thick},
}

% --- utilidades -------------------------------------------------------------
% Identificadores y rutas con guion bajo, sin escaparlos a mano.
\newcommand{\id}[1]{\texttt{\detokenize{#1}}}
% Rutas largas en celdas angostas: \path corta en / y en guiones.
\newcommand{\ruta}{\path}
% Dato que el grupo debe completar antes de entregar. Buscar "PENDIENTE".
\newcommand{\pendiente}[1]{\textcolor{cPend}{\textbf{[PENDIENTE: #1]}}}
\newcolumntype{L}[1]{>{\raggedright\arraybackslash}p{#1}}
\newcolumntype{C}[1]{>{\centering\arraybackslash}p{#1}}
\newcolumntype{Y}{>{\raggedright\arraybackslash}X}

\newsavebox{\notabox}
\newenvironment{nota}
  {\par\medskip\begin{lrbox}{\notabox}%
   \begin{minipage}{\dimexpr0.95\linewidth-14pt\relax}\small}
  {\end{minipage}\end{lrbox}%
   \begin{center}\setlength{\fboxsep}{7pt}\colorbox{cFondo}{\usebox{\notabox}}\end{center}\medskip}

% ---------------------------------------------------------------------------
\begin{document}

\begin{titlepage}
\centering
{\large\bfseries Instituto Tecnológico de Costa Rica\par}
\vspace{0.3cm}
{\large Escuela de Ingeniería en Computadores\par}
\vspace{0.3cm}
{\large CE-1113 Sistemas Empotrados\par}
\vspace{3.2cm}
{\LARGE\bfseries Documento de Diseño (DI)\par}
\vspace{0.6cm}
{\Large Proyecto I --- Sistema embebido a la medida para un robot aspiradora
autónomo con reproducción de audio y control remoto\par}
\vspace{3.2cm}
{\large\bfseries Integrantes\par}
\vspace{0.2cm}
{\large \pendiente{nombres y carnés de los cuatro integrantes}\par}
\vspace{1.6cm}
{\large\bfseries Profesor\par}
\vspace{0.2cm}
{\large Dr.-Ing. Jeferson González Gómez\par}
\vfill
{\large Octubre de 2026\par}
\end{titlepage}

\tableofcontents
\newpage

% ===========================================================================
\section{Introducción}

Este documento evidencia el atributo de \textbf{Diseño (DI)} en el Proyecto I
del curso. Recorre el proceso completo: qué se necesitaba y bajo qué
restricciones (DI1), qué alternativas se valoraron y con qué criterios (DI2),
cómo quedó diseñada la alternativa elegida (DI3) y cómo se comprobó que el
resultado cumple los requerimientos y las consideraciones de seguridad, costo
e impacto (DI4).

El producto es un robot aspiradora de tracción diferencial sobre una
Raspberry Pi 4 con una imagen de Linux mínima construida con Yocto (Poky,
rama \id{scarthgap}). Navega de forma autónoma evadiendo obstáculos, reproduce
MP3 y se opera desde un panel web que muestra la telemetría y el mapa de
recorrido en tiempo real. Todo acceso al hardware pasa por una biblioteca
dinámica propia, \id{librobot.so}, compilada de forma cruzada.

El detalle técnico ---cálculos, listas de materiales, procedimientos de
verificación y evidencias--- está en la carpeta \id{docs/} del repositorio.
Aquí se resume y se cita el archivo correspondiente.

% ===========================================================================
\section{DI1 --- Necesidades y requerimientos}

\subsection{El problema}

Un robot de limpieza coordina en tiempo real subsistemas que compiten por los
mismos recursos: sensado de proximidad, tracción, audio, red e interfaz de
usuario. El problema es complejo por tres razones que no se pueden resolver
por separado:

\begin{itemize}
  \item \textbf{Recursos acotados.} El enunciado fija un presupuesto de
        200\,MB de sistema de archivos raíz y 15\,s de arranque. Los GPIO
        entregan 16\,mA por pin y 50\,mA en total, y no toleran más de
        3.3\,V.
  \item \textbf{Dominios eléctricos incompatibles.} La lógica trabaja a
        3.3\,V y miliamperios; los motores, a 7.4\,V con picos de varias
        veces su corriente nominal. Deben convivir sin compartir tierra.
  \item \textbf{Concurrencia.} Navegar, medir, reproducir audio y atender la
        interfaz ocurren a la vez, y ninguna tarea puede degradar la
        respuesta ante un obstáculo.
\end{itemize}

\subsection{Partes interesadas y sus necesidades}

\begin{center}
\begin{tabularx}{\linewidth}{L{3.1cm}Y}
\toprule
\textbf{Parte interesada} & \textbf{Necesidad} \\
\midrule
Usuario final & Un robot que limpie sin supervisión, que se pueda tomar a
mano desde el celular y cuyo estado se entienda de un vistazo (LEDs, sonidos,
mapa). \\
Curso y profesor & Cumplimiento verificable del enunciado, y que el equipo
prestado (Raspberry Pi 4 y su kit) regrese sin daño. \\
Grupo de trabajo & Un diseño que cuatro personas puedan construir en seis
semanas, con el inventario disponible y un costo que el grupo pueda asumir. \\
Terceros en el entorno & Que el robot no represente un riesgo: ni eléctrico,
ni de incendio, ni por golpear a personas u objetos, ni por exponer la red. \\
\bottomrule
\end{tabularx}
\end{center}

\subsection{Requerimientos técnicos}

Los requerimientos obligatorios del enunciado se numeran para trazarlos
después en la validación (sección~\ref{sec:matriz}).

\begin{longtable}{L{1.2cm}L{3.6cm}L{10.2cm}}
\toprule
\textbf{Id.} & \textbf{Área} & \textbf{Requerimiento} \\
\midrule
\endhead
R1 & Navegación & Modo autónomo con al menos un algoritmo de cobertura
reactiva. Ante un obstáculo: detenerse, retroceder y cambiar de dirección. \\
R2 & Sensado & Al menos dos sensores de proximidad para detección frontal y
lateral, leídos y procesados en tiempo real. \\
R3 & Tracción & Dos motores DC con control diferencial (avance, retroceso,
giros, detención) y velocidad por PWM para giros de radio variable. \\
R4 & Indicadores & Cuatro LEDs: modo autónomo, modo manual, obstáculo y
sistema encendido. \\
R5 & Audio & Reproducir MP3 locales de forma concurrente con la navegación. \\
R6 & Audio & Desde la interfaz: elegir canción, reproducir, pausar, detener y
ajustar el volumen. \\
R7 & Audio & Salida directa desde el robot, con parlante a bordo. \\
R8 & Audio & Sonidos de notificación en cuatro eventos: inicio del sistema,
inicio del modo autónomo, obstáculo y cambio a manual. \\
R9 & Control remoto & Conmutar entre modo autónomo y manual; en manual, los
controles direccionales mandan los motores. \\
R10 & Control remoto & Panel con modo activo, controles, sensores en tiempo
real, control de audio, estado de los LEDs y mapa. \\
R11 & Mapa & Grilla 2D incremental (visitada, obstáculo, desconocida) a
partir de los sensores y la odometría, visible en tiempo real. \\
R12 & Seguridad & Pantalla de inicio de sesión con al menos un usuario y un
protocolo de seguridad mínimo. \\
R13 & Arranque & Inicio automático con una unidad \id{systemd} propia,
\id{Restart=on-failure} y sin interfaz gráfica local. \\
R14 & Biblioteca & Biblioteca dinámica propia que encapsule motores, sensores,
LEDs y audio; el servidor la usa de forma exclusiva. \\
R15 & Desarrollo & Todo el software se compila de forma cruzada en el host. \\
R16 & Sistema operativo & Imagen mínima con Yocto, solo con los paquetes
necesarios y cada uno justificado. \\
R17 & Construcción & CMake o Autotools, generando el paquete estándar. \\
R18 & Receta & Al menos una receta \id{.bb} propia; la imagen se reproduce
con \id{bitbake <imagen>} sin pasos manuales. \\
R19 & Eficiencia & Medir y reportar rootfs, tiempo de arranque, RAM y CPU en
operación normal. Referencia: rootfs $\leq$ 200\,MB y arranque $\leq$ 15\,s. \\
\bottomrule
\end{longtable}

Los tres requerimientos opcionales (detección de desnivel, notificación de
fin de ciclo y playlist persistente) se dejaron fuera del alcance para
concentrar el tiempo en el núcleo obligatorio.

\subsection{Salud y seguridad pública}

\begin{center}
\begin{tabularx}{\linewidth}{L{3.6cm}Y}
\toprule
\textbf{Riesgo} & \textbf{Requerimiento que se deriva} \\
\midrule
Incendio o fuga térmica de las celdas de Li-Ion & BMS con protección de
sobrecarga, sobredescarga, cortocircuito y balanceo; fusible en serie con el
pack; nunca alimentar la Raspberry directo de la celda. \\
Daño al equipo prestado & Aislamiento óptico entre la lógica y la etapa de
potencia, con tierras separadas; ninguna señal de 5\,V llega a un GPIO. \\
Golpes a personas u objetos & Freno automático ante un obstáculo también en
modo manual; masa menor a 1.5\,kg; los motores quedan frenados en reposo y
durante el arranque. \\
Cables en las llantas & Ruteo fijo del cableado, sin tramos sueltos a menos
de 1\,cm de una llanta. \\
Acceso no autorizado por la red & Inicio de sesión obligatorio; el demonio
de GPIO solo escucha en la interfaz local; sin acceso remoto sin contraseña
en la imagen de entrega. \\
Nivel sonoro & Volumen ajustable desde la interfaz y potencia máxima acotada
por el divisor de entrada del amplificador. \\
\bottomrule
\end{tabularx}
\end{center}

\subsection{Costos}

La Raspberry Pi 4 y su kit los presta el curso, de modo que el costo para el
grupo es el del modelo físico. La tabla es la estimación con que se planeó la
compra, por bloques.

\begin{center}
\begin{tabular}{L{6.6cm}C{3.2cm}}
\toprule
\textbf{Bloque} & \textbf{Estimación (USD)} \\
\midrule
Alimentación portátil (pack 2S, BMS, reguladores, fusible, interruptor) & 25--45 \\
Circuitería y sensores (optoacopladores, servo, pasivos, placas) & 9--17 \\
Etapa de audio (PAM8403, parlante, filtro) & 5--12 \\
Chasis y acabado & 10--25 \\
\midrule
\textbf{Total del robot} & \textbf{49--99} \\
\bottomrule
\end{tabular}
\end{center}

\pendiente{sustituir la estimación por el costo real, con las facturas de la
compra}

El requerimiento que se fijó fue mantenerse dentro de ese intervalo y
aprovechar primero lo que el grupo ya tenía: un sensor ultrasónico, un
MPU-6050 y un driver L298N.

\subsection{Impacto ambiental}

\begin{itemize}
  \item \textbf{Baterías.} Las celdas de Li-Ion son el componente de mayor
        impacto. Se exige que sean recargables y desmontables, y que al final
        de su vida se entreguen a un gestor de residuos electrónicos.
  \item \textbf{Residuos electrónicos.} El montaje debe permitir reutilizar
        los módulos: conectores desconectables y nada pegado de forma
        permanente a la Raspberry Pi, que se devuelve al curso.
  \item \textbf{Consumo.} Menos consumo es más autonomía con el mismo pack y
        menos ciclos de carga: se prefieren reguladores conmutados y
        amplificación clase D.
  \item \textbf{Compras.} Reutilizar el inventario del grupo antes que
        comprar componentes nuevos.
\end{itemize}

\subsection{Recursos disponibles y restricciones}

\begin{center}
\begin{tabularx}{\linewidth}{L{3.6cm}Y}
\toprule
\textbf{Recurso} & \textbf{Disponibilidad} \\
\midrule
Plataforma & Raspberry Pi 4 y kit, prestados por el curso. \\
Inventario del grupo & Un HC-SR04, un MPU-6050 (módulo GY-521) y un driver
L298N de doble puente. \\
Personas y tiempo & Cuatro integrantes, del 25 de agosto al 6 de octubre de
2026. \\
Host de construcción & Laptop de 4 núcleos y 7\,GB de RAM, por debajo de lo
recomendado para Yocto. \\
Software de partida & Un proyecto de referencia con licencia MIT, del que
derivan parte de la capa y del código (atribuido en \id{NOTICE.md}). \\
Herramientas & Multímetro, cautín, y el Toolchain-SDK generado por el propio
proyecto. \\
\bottomrule
\end{tabularx}
\end{center}

De ahí salen dos restricciones que marcaron el diseño: el hardware llegó
tarde respecto del software, así que hacía falta poder desarrollar y probar
sin él; y el host de construcción era el cuello de botella de cada iteración
de la imagen.

% ===========================================================================
\section{DI2 --- Valoración de alternativas}

\subsection{Método}

Las decisiones de mayor peso se evaluaron con una matriz ponderada. Cada
alternativa recibe una nota de 1 (peor) a 5 (mejor) en cuatro factores:

\begin{center}
\begin{tabularx}{\linewidth}{L{2.4cm}C{1.4cm}Y}
\toprule
\textbf{Factor} & \textbf{Peso} & \textbf{Qué mide} \\
\midrule
Técnico & 40\,\% & Cumplimiento del requerimiento, desempeño, uso de GPIO y
riesgo de integración. \\
Económico & 25\,\% & Costo de compra y disponibilidad local; cuánto del
inventario se aprovecha. \\
Ambiental & 15\,\% & Cantidad de componentes, consumo y reutilización. \\
Social & 20\,\% & Seguridad de las personas y del equipo prestado. \\
\bottomrule
\end{tabularx}
\end{center}

El factor técnico pesa más porque una alternativa que no cumple el
requerimiento no es una alternativa. El social pesa más que el ambiental
porque el enunciado hace al grupo responsable del equipo prestado.

\subsection{Detección de obstáculos}

\begin{center}
\begin{tabularx}{\linewidth}{Y C{1.25cm}C{1.4cm}C{1.4cm}C{1.1cm}C{1.2cm}}
\toprule
\textbf{Alternativa} & \textbf{Téc.} & \textbf{Econ.} & \textbf{Amb.} &
\textbf{Soc.} & \textbf{Total} \\
\midrule
Tres HC-SR04 fijos & 3 & 3 & 3 & 4 & 3.20 \\
\textbf{Un HC-SR04 sobre un servo (radar)} & 4 & 4 & 4 & 3 & \textbf{3.80} \\
Sensores infrarrojos & 2 & 4 & 3 & 2 & 2.65 \\
\bottomrule
\end{tabularx}
\end{center}

Los tres sensores fijos dan una lectura frontal rápida ($\approx$ 16\,Hz),
pero solo tres direcciones, ocupan seis GPIO, obligan a secuenciar los
disparos para evitar ecos cruzados y exigían comprar dos sensores más. Los
infrarrojos económicos tienen poco alcance, dependen del color de la
superficie y de la luz ambiente, y no entregan una distancia útil para el
mapa.

El radar mide siete direcciones con tres GPIO y el sensor que el grupo ya
tenía. Su debilidad es que mira al frente una vez cada 1.2\,s; por eso su
nota social es menor. Esa debilidad se atacó en el diseño con la estimación
del tiempo antes de chocar (sección~\ref{sec:creativo}).

\subsection{Alimentación}

\begin{center}
\begin{tabularx}{\linewidth}{Y C{1.25cm}C{1.4cm}C{1.4cm}C{1.1cm}C{1.2cm}}
\toprule
\textbf{Alternativa} & \textbf{Téc.} & \textbf{Econ.} & \textbf{Amb.} &
\textbf{Soc.} & \textbf{Total} \\
\midrule
Una celda (1S) con elevador a 5\,V & 2 & 4 & 4 & 2 & 2.80 \\
\textbf{Pack 2S con BMS y reductores} & 5 & 3 & 3 & 5 & \textbf{4.20} \\
Powerbank USB y batería aparte para motores & 3 & 3 & 2 & 4 & 3.05 \\
\bottomrule
\end{tabularx}
\end{center}

Con una celda, entregar 15\,W a 5\,V pide cerca de 4.5\,A de entrada, y más
de 5\,A con la celda descargada: los elevadores MT3608 y XL6009 no lo
sostienen, y una caída de tensión corrompe la microSD. El powerbank resuelve
la lógica pero obliga a un segundo pack para los motores y deja de ser una
sola fuente regulada por separado. Con 2S, un reductor trabaja con holgura en
todo el rango de descarga (6.0--8.4\,V) y los motores se alimentan directo.

\subsection{Salida de audio}

\begin{center}
\begin{tabularx}{\linewidth}{Y C{1.25cm}C{1.4cm}C{1.4cm}C{1.1cm}C{1.2cm}}
\toprule
\textbf{Alternativa} & \textbf{Téc.} & \textbf{Econ.} & \textbf{Amb.} &
\textbf{Soc.} & \textbf{Total} \\
\midrule
Jack de 3.5\,mm con LM386 & 3 & 4 & 3 & 3 & 3.25 \\
\textbf{PWM por GPIO 18 con PAM8403} & 4 & 5 & 4 & 4 & \textbf{4.25} \\
DAC I2S (MAX98357A) & 4 & 3 & 4 & 4 & 3.75 \\
\bottomrule
\end{tabularx}
\end{center}

El LM386 es clase AB, rinde cerca del 50\,\% y pide ocho componentes
discretos para dar 0.5\,W. El DAC I2S da el mejor audio, pero es más caro,
más difícil de conseguir localmente y cambia la configuración de ALSA. El
PAM8403 es un módulo armado, clase D ($\approx$ 85\,\%), que entrega 1.3\,W
sobre 8\,$\Omega$ con un filtro de tres pasivos.

\subsection{Odometría}

\begin{center}
\begin{tabularx}{\linewidth}{Y C{1.25cm}C{1.4cm}C{1.4cm}C{1.1cm}C{1.2cm}}
\toprule
\textbf{Alternativa} & \textbf{Téc.} & \textbf{Econ.} & \textbf{Amb.} &
\textbf{Soc.} & \textbf{Total} \\
\midrule
Solo el modelo de los motores & 2 & 5 & 5 & 3 & 3.40 \\
Encoders ópticos de rueda & 4 & 3 & 3 & 3 & 3.40 \\
\textbf{MPU-6050 fusionado con el modelo} & 4 & 5 & 5 & 4 & \textbf{4.40} \\
\bottomrule
\end{tabularx}
\end{center}

Estimar la posición solo con el tiempo y el ciclo de trabajo deja derivar el
rumbo en cuanto las dos llantas giran distinto, y el rumbo es lo que más
deforma el mapa. Los encoders miden bien la distancia, pero había que
comprarlos y montarlos en motores que no traen eje para ello. El MPU-6050 ya
estaba en el inventario: su giroscopio da el rumbo y su acelerómetro la
dinámica de la velocidad, que además habilita el tiempo antes de chocar.

\subsection{Otras decisiones}

\begin{longtable}{L{2.6cm}L{3.8cm}L{8.6cm}}
\toprule
\textbf{Decisión} & \textbf{Elegida} & \textbf{Alternativas y motivo} \\
\midrule
\endhead
Aislamiento & Cinco PC817 & El DRV8871 acepta 3.3\,V pero comparte tierra:
no aísla y se descartó. El ADuM1401 aísla cuatro canales en un chip, pero es
más caro y difícil de conseguir. \\
Acceso a GPIO & \id{pigpio} & \id{libgpiod} y WiringPi no generan PWM ni
pulsos de servo con temporización estable; \id{pigpio} lo hace por DMA, mide
el eco en microsegundos y da el bus I2C por el mismo demonio. \\
Servidor web & \id{libmicrohttpd} & nginx o lighttpd son un proceso aparte
con CGI que mantener; un servidor en Python suma unos 40\,MB al rootfs. \\
Decodificador MP3 & \id{mpg123} & GStreamer trae decenas de MB de plugins
para el mismo resultado. \\
Imagen base & \id{core-image-minimal} & \id{core-image-base} y
\id{full-cmdline} arrastran utilidades y servicios que el robot no usa. \\
Chasis & Circular & El rectangular es más fácil de cortar, pero se traba en
las esquinas al girar; el circular rota sobre su eje sin barrer área extra. \\
Algoritmo de cobertura & Aleatorio con rebote & Espiral y zig-zag dependen
de una odometría precisa; el rebote tolera el error acumulado y saca partido
de las siete direcciones del radar. \\
\bottomrule
\end{longtable}

% ===========================================================================
\section{DI3 --- Diseño de la alternativa seleccionada}

\subsection{Arquitectura de hardware}

Un pack 2S alimenta dos rieles regulados por separado. Las cinco señales que
cruzan del dominio lógico al de potencia ---las cuatro entradas del puente H
y la del servo--- lo hacen por optoacopladores, y las dos tierras no se
tocan en ningún punto.

\begin{figure}[h]
\centering
\resizebox{\linewidth}{!}{%
\begin{tikzpicture}
  \node[blkn] (bat)  at (0,0)       {Pack 2S 18650\\7.4\,V};
  \node[blkn] (bms)  at (3.1,0)     {BMS 2S\\fusible 5\,A};

  \node[blkl] (buck) at (6.5,1.7)   {Buck MP2307\\5\,V / 3\,A};
  \node[blkl] (pi)   at (10.0,1.7)  {Raspberry Pi 4\\\texttt{robot-image}};
  \node[blkl] (per)  at (14.0,1.7)  {HC-SR04 $\cdot$ MPU-6050\\4 LEDs $\cdot$ PAM8403};

  \node[blkn] (opto) at (10.0,0)    {5 $\times$ PC817};

  \node[blkp] (l298) at (10.0,-1.7) {L298N\\(\texttt{ENA}/\texttt{ENB} con jumper)};
  \node[blkp] (mot)  at (14.0,-1.7) {2 motores DC};
  \node[blkp] (bk2)  at (6.5,-3.4)  {Buck MP2307\\5\,V (servo)};
  \node[blkp] (srv)  at (14.0,-3.4) {Servo del radar};

  \draw[ln] (bat) -- (bms);
  \draw[ln] (bms.east) -- ++(0.5,0) |- (buck.west);
  \draw[ln] (bms.east) -- ++(0.5,0) |- (l298.west);
  \draw[ln] (bms.east) -- ++(0.5,0) |- (bk2.west);
  \draw[ln] (buck) -- (pi);
  \draw[ln] (pi) -- (per);
  \draw[ln,dashed] (pi) -- (opto);
  \draw[ln,dashed] (opto) -- node[right,font=\scriptsize]{\texttt{IN1}--\texttt{IN4}} (l298);
  \draw[ln,dashed] (opto.east) -- (16.2,0) |- node[pos=0.25,right,font=\scriptsize]{pulsos} (srv.east);
  \draw[ln] (l298) -- (mot);
  \draw[ln] (bk2) -- (srv);

  \begin{scope}[on background layer]
    \node[draw=cLog,dashed,rounded corners,fit=(buck)(pi)(per),inner sep=7pt,
          label={[font=\scriptsize\bfseries,text=cLog]above:Dominio lógico --- GND\_LOG}] {};
    \node[draw=cPot,dashed,rounded corners,fit=(l298)(mot)(bk2)(srv),inner sep=7pt,
          label={[font=\scriptsize\bfseries,text=cPot]below:Dominio de potencia --- GND\_POT}] {};
  \end{scope}
\end{tikzpicture}}
\caption{Arquitectura de hardware. Las líneas punteadas son señales que
cruzan la barrera óptica. Los esquemas eléctricos completos están en
\texttt{hardware-diagramas.pdf}.}
\end{figure}

\paragraph{Alimentación.} Dos celdas 18650 en serie con BMS de balanceo y un
fusible de 5\,A. El riel lógico sale de un reductor MP2307 ajustado a 5.0\,V
(3\,A, $\approx$ 93\,\% de eficiencia) y alimenta la Raspberry Pi por USB-C,
el sensor, el MPU-6050, los LEDs y el amplificador: 2.0\,A típicos y 3.0\,A
de pico. El riel de potencia lleva los 7.4\,V directo al L298N. El servo
tiene su propio reductor de 5\,V en el dominio de potencia, porque sus picos
de 700\,mA no caben ni en el riel lógico ni en el regulador interno del
L298N. Autonomía estimada: 60 minutos. Detalle en
\id{docs/hardware-alimentacion.md}.

\paragraph{Aislamiento.} Un PC817 por señal. Resistencia de entrada de
390\,$\Omega$, que fija 5\,mA por canal: con cuatro canales son 20\,mA, y
junto con los LEDs y el servo el conector entrega 39\,mA de los 50\,mA que
admite. El L298N lleva los jumpers \id{ENA} y \id{ENB} puestos y se controla
solo con \id{IN1}--\id{IN4}, lo que ahorra dos optoacopladores y 10\,mA.
Detalle en \id{docs/hardware-aislamiento.md}.

\paragraph{Sensado.} Un HC-SR04 sobre un servo de 180° en el borde frontal,
con el eje 10\,cm por delante del centro. El pin \id{ECHO} entrega 5\,V y
pasa por un divisor de 1\,k$\Omega$ y 2\,k$\Omega$ (3.33\,V) antes del GPIO.
El MPU-6050 va entre las dos llantas, alimentado a 3.3\,V para que sus
resistencias de pull-up nunca pongan 5\,V en el bus I2C. Detalle en
\id{docs/hardware-sensores.md}.

\paragraph{Indicadores y audio.} Cuatro LEDs de colores distintos a 5\,mA,
con la resistencia calculada para cada color. El audio sale por PWM en el
GPIO 18 (overlay \id{audremap}), pasa por un filtro RC con corte en 19\,kHz y
entra al PAM8403, alimentado del riel lógico para que el ruido de los motores
no llegue al parlante.

\paragraph{Modelo físico.} Chasis circular de 22--25\,cm en dos niveles, con
la batería al centro y abajo para bajar el centro de gravedad, el eje de
tracción ligeramente adelante del centro y una rueda loca atrás. Detalle en
\id{docs/hardware-chasis.md}.

\subsection{Arquitectura de software}

\begin{figure}[h]
\centering
\resizebox{\linewidth}{!}{%
\begin{tikzpicture}[node distance=5mm]
  \tikzset{capa/.style={blk,minimum width=13.2cm},
           media/.style={blk,minimum width=6.4cm}}
  \node[capa,blkn] (web) {Navegador del celular o la PC: \texttt{login.html} y \texttt{dashboard.html}};
  \node[capa,blkl,below=of web] (srv) {\textbf{\texttt{robot-server}}: API REST $\cdot$ autenticación $\cdot$ hilo de navegación (10\,Hz) $\cdot$ mapa de 31 $\times$ 31};
  \node[capa,blkl,below=of srv] (lib) {\textbf{\texttt{librobot.so.1}}: motores $\cdot$ servo $\cdot$ radar (hilo) $\cdot$ MPU-6050 $\cdot$ odometría (hilo, 50\,Hz) $\cdot$ LEDs $\cdot$ audio (hilo)};
  \node[media,blkn,below=of lib.south west,anchor=north west] (pig) {\texttt{pigpiod}: GPIO, PWM y pulsos por DMA, I2C};
  \node[media,blkn,below=of lib.south east,anchor=north east] (snd) {\texttt{mpg123} y ALSA};
  \node[capa,blkn,below=24mm of lib] (ker) {Linux 6.6 y \texttt{systemd} --- imagen \texttt{robot-image} (Yocto \texttt{scarthgap})};
  \node[capa,blkp,below=of ker] (hw) {Raspberry Pi 4 $\cdot$ L298N $\cdot$ servo $\cdot$ HC-SR04 $\cdot$ MPU-6050 $\cdot$ LEDs $\cdot$ PAM8403};

  \draw[ln] (web) -- node[right,font=\scriptsize]{HTTP en el puerto 8080, JSON, sondeo cada 500\,ms} (srv);
  \draw[ln] (srv) -- node[right,font=\scriptsize]{única vía de acceso al hardware} (lib);
  \draw[ln] (lib.south -| pig) -- (pig);
  \draw[ln] (lib.south -| snd) -- (snd);
  \draw[ln] (pig) -- (pig |- ker.north);
  \draw[ln] (snd) -- (snd |- ker.north);
  \draw[ln] (ker) -- (hw);
\end{tikzpicture}}
\caption{Arquitectura de software por capas.}
\end{figure}

\paragraph{Sistema operativo.} La imagen \id{robot-image} parte de
\id{core-image-minimal} y agrega solo lo que el robot usa; cada paquete está
justificado en \id{docs/paquetes.md}. La capa propia \id{meta-robot} aporta
las recetas de \id{librobot}, \id{robot-server}, \id{pigpio}, la
configuración de ALSA y de WiFi, los ajustes de \id{config.txt} y la tabla
de particiones. La música va en una tercera partición, de modo que ninguna
supera los 200\,MB.

\paragraph{Biblioteca.} \id{librobot} es la única vía de acceso al hardware.
Una llamada, \id{robot_init()}, abre la sesión e inicia dos hilos propios: la
odometría a 50\,Hz y el barrido del radar. El audio corre en un tercer hilo.
La referencia completa está en \id{docs/api-librobot.md}.

\paragraph{Servidor.} \id{robot-server} sirve la interfaz, expone la API y
corre el hilo de navegación a 10\,Hz. No espera al radar: lee la última
lectura de cada ángulo cuando la necesita. Las contraseñas se guardan como
resumen SHA-256 y la sesión viaja en una cookie \id{HttpOnly} que caduca a
los 10 minutos de inactividad.

\paragraph{Arranque.} Dos unidades \id{systemd}, habilitadas desde las
recetas. \id{robot-server} declara \id{Requires=pigpiod.service} y
\id{Restart=on-failure} con un límite de cinco reintentos por minuto.

\subsection{Navegación}

El robot avanza hasta que detecta un obstáculo por cualquiera de dos vías:
el tiempo antes de chocar baja de 1.2\,s, o alguna lectura del cono frontal
(60°, 90°, 120°) baja de 20\,cm. Entonces se detiene, suena el aviso y
enciende el LED rojo. En modo autónomo sigue la evasión:

\begin{enumerate}
  \item retrocede medio segundo;
  \item espera un barrido completo del radar, ya quieto;
  \item elige el rumbo con más espacio entre las seis direcciones que no son
        el frente; las que quedan a menos de 25\,cm de la mejor compiten al
        azar, que es la parte aleatoria del rebote;
  \item gira en lazo cerrado sobre el rumbo del giroscopio hasta cubrir el
        ángulo elegido, no durante un tiempo fijo.
\end{enumerate}

En modo manual la misma detección actúa como freno de seguridad.

\subsection{Elementos propios del diseño}
\label{sec:creativo}

\begin{itemize}
  \item \textbf{Tiempo antes de chocar.} Cada lectura frontal se congela con
        el avance acumulado en ese instante. Entre dos lecturas, la distancia
        se proyecta con lo que el robot avanzó y se divide entre la velocidad
        actual: $t = (d - \Delta\text{avance}) / v$. La cuenta regresiva
        sigue bajando aunque el sensor esté mirando a un lado, y compensa la
        baja frecuencia de la lectura frontal del radar.
  \item \textbf{Fusión de la velocidad.} Un filtro complementario deja pasar
        la dinámica rápida del acelerómetro y toma del modelo de los motores
        la tendencia lenta ($\tau = 1$\,s); con los motores detenidos la
        velocidad se fuerza a cero. Así el sesgo del acelerómetro no se
        acumula.
  \item \textbf{El canal del servo no invierte.} Los cuatro canales del
        puente H van en emisor común, que invierte la señal, y la inversión
        se compensa en un único punto de \id{lib_motors.c}. Para el servo eso
        no sirve: en reposo el GPIO queda en bajo, y un canal inversor le
        entregaría un nivel alto permanente que lo llevaría contra su tope.
        Su canal se armó en seguidor de emisor, que no invierte.
  \item \textbf{El reposo es seguro.} Con los GPIO en bajo, las cuatro
        entradas del L298N quedan en alto por sus resistencias: con el puente
        habilitado eso es freno, no movimiento.
  \item \textbf{Desarrollo sin hardware.} Un simulador reemplaza a
        \id{pigpiod} por un robot virtual en una sala de 4 $\times$ 4\,m, sin
        tocar una línea de la biblioteca. Una máquina \id{qemuarm64-robot}
        con el mismo \emph{tune} Cortex-A72 arranca la imagen completa con
        los mismos binarios de espacio de usuario que corren en la Pi.
  \item \textbf{Imagen reproducible sin secretos.} Si falta el archivo de
        credenciales de WiFi, la receta usa la plantilla y avisa: un clon
        limpio produce la imagen sin pasos manuales, como exige R18.
\end{itemize}

\subsection{Cómo responde el diseño a seguridad, costo e impacto}

\begin{center}
\begin{tabularx}{\linewidth}{L{2.6cm}Y}
\toprule
\textbf{Aspecto} & \textbf{Respuesta del diseño} \\
\midrule
Seguridad & BMS, fusible, dos rieles, aislamiento óptico con tierras
separadas, divisor en \id{ECHO}, MPU-6050 a 3.3\,V, freno en reposo y freno
automático en modo manual, \id{pigpiod} con \id{-l} (solo interfaz local),
inicio de sesión obligatorio. \\
Costo & Se reutilizó el sensor, el MPU-6050 y el L298N del grupo. El radar
evitó comprar dos sensores; el MPU-6050 evitó los encoders; el PAM8403
reemplazó ocho componentes discretos. \\
Ambiental & Reguladores conmutados y amplificador clase D; módulos con
conectores desconectables y reutilizables; celdas recargables y
desmontables. \\
\bottomrule
\end{tabularx}
\end{center}

% ===========================================================================
\section{DI4 --- Validación del diseño final}

\subsection{Estrategia}

La validación se hizo en tres niveles, de menor a mayor costo de cada
iteración:

\begin{enumerate}
  \item \textbf{Simulador en el host.} \id{sim/prueba_librobot} ejercita
        cada módulo de la biblioteca contra el robot virtual y compara con la
        verdad del mundo simulado: 47 comprobaciones.
  \item \textbf{Emulación ARM.} El mismo programa, compilado con el
        Toolchain-SDK, corre bajo \id{qemu-aarch64}; y la imagen completa
        arranca en la máquina \id{qemuarm64-robot}.
  \item \textbf{Target.} La imagen en la Raspberry Pi 4 con el robot armado,
        siguiendo los procedimientos de verificación de \id{docs/}.
\end{enumerate}

\subsection{Matriz de cumplimiento}
\label{sec:matriz}

\begin{longtable}{L{0.8cm}L{6.3cm}L{1.9cm}L{5.4cm}}
\toprule
\textbf{Id.} & \textbf{Cómo se validó} & \textbf{Resultado} & \textbf{Evidencia} \\
\midrule
\endhead
R1 & Dos minutos en modo autónomo sobre el simulador: 18 evasiones y ningún
choque. Prueba de campo en un área delimitada. & Cumple &
\ruta{docs/navegacion-radar.md} \\
R2 & Barrido completo contra el mundo simulado: los siete ángulos a menos de
2\,cm de la distancia real, a 5 lecturas por segundo. & Cumple &
\ruta{sim/prueba_librobot.c} \\
R3 & Avance, retroceso, giros y saturación verificados en el simulador;
sentido de giro confirmado con el robot elevado. & Cumple &
\ruta{docs/evidencias/ejecucion-cruzada-qemu.md} \\
R4 & Cada LED cambia de estado al ordenarlo y sigue al modo activo. & Cumple
& \ruta{docs/evidencias/servidor-web.md} \\
R5 & Reproducción en un hilo propio mientras el hilo de navegación corre. &
Cumple & \ruta{lib/lib_audio.c} \\
R6 & Listar, reproducir, pausar, detener y volumen por HTTP. & Cumple &
\ruta{sim/prueba_api.sh} \\
R7 & Tono de prueba y sonido de evento por el parlante del robot. & Cumple &
\ruta{docs/hardware-sensores.md} \\
R8 & Los cuatro eventos disparan su sonido; el aviso pausa la música y la
reanuda. & Cumple & \ruta{docs/api-librobot.md} \\
R9 & Cambio de modo y movimiento manual por HTTP; los controles manuales se
bloquean en autónomo. & Cumple & \ruta{docs/evidencias/servidor-web.md} \\
R10 & \id{/api/status} entrega modo, sensores, LEDs, audio y mapa en una sola
consulta, cada 500\,ms. & Cumple & \ruta{docs/evidencias/servidor-web.md} \\
R11 & La grilla crece en tiempo real: de 79 a 85 celdas visitadas en 6\,s. &
Cumple & \ruta{docs/evidencias/servidor-web.md} \\
R12 & Credenciales inválidas y consultas sin sesión devuelven 401. & Cumple
& \ruta{docs/evidencias/servidor-web.md} \\
R13 & Tras un \id{kill -9}, el servicio vuelve con otro PID y
\id{NRestarts=1}. & Cumple & \ruta{docs/evidencias/systemd-reinicio.md} \\
R14 & El servidor declara \id{NEEDED librobot.so.1} y no toca GPIO. & Cumple
& \ruta{docs/evidencias/binarios-target.md} \\
R15 & Los logs muestran \id{aarch64-poky-linux-gcc}; los binarios son ELF
ARM aarch64 y no ejecutan en el host. & Cumple & \ruta{docs/evidencias/} \\
R16 & Imagen sobre \id{core-image-minimal}; cada paquete agregado, listado y
justificado. & Cumple & \ruta{docs/paquetes.md} \\
R17 & Ambas recetas heredan \id{cmake}; la biblioteca se empaqueta en
\id{librobot} y \id{librobot-dev}. & Cumple & \ruta{meta-robot/} \\
R18 & \id{bitbake robot-image} produce la imagen desde un clon limpio. &
Cumple & \ruta{docs/yocto-setup.md} \\
R19 & Ver la sección~\ref{sec:metricas}. & Ver tabla & \ruta{docs/metricas.md} \\
\bottomrule
\end{longtable}

\pendiente{R3: confirmar el estado final de la velocidad variable; ver el
hallazgo 3 de la sección~\ref{sec:hallazgos}}

\subsection{Validación de la seguridad}

Antes de conectar la Raspberry Pi se ejecutaron, en orden, los procedimientos
de cinco pasos de \id{docs/hardware-aislamiento.md} y
\id{docs/hardware-alimentacion.md}. Los criterios de aceptación:

\begin{center}
\begin{tabularx}{\linewidth}{Y L{5.6cm}}
\toprule
\textbf{Verificación} & \textbf{Criterio} \\
\midrule
Continuidad entre \id{GND_LOG} y \id{GND_POT}, todo desenergizado &
Circuito abierto ($>$ 1\,M$\Omega$) \\
Diferencia de tensión entre las dos celdas & Menor a 0.05\,V \\
Salida del reductor de 5\,V en vacío y con 2\,A de carga &
5.00 $\pm$ 0.10\,V y $\geq$ 4.90\,V \\
Salida del divisor de \id{ECHO} al disparar el sensor & 3.3\,V como máximo \\
Entradas del L298N con los optoacopladores en reposo &
$\approx$ 5\,V: motores frenados \\
Salida del canal del servo con el GPIO en bajo & $\approx$ 0\,V \\
\id{vcgencmd get_throttled} navegando con audio & \id{0x0} \\
\bottomrule
\end{tabularx}
\end{center}

\pendiente{anotar los valores medidos en cada verificación}

\subsection{Validación de la eficiencia}
\label{sec:metricas}

\begin{center}
\begin{tabularx}{\linewidth}{L{2.9cm}C{1.9cm}C{3.1cm}Y}
\toprule
\textbf{Métrica} & \textbf{Referencia} & \textbf{Medido} & \textbf{Método} \\
\midrule
Rootfs & $\leq$ 200\,MB & 129\,MB & \id{du} sobre el rootfs de la imagen
construida. La partición raíz es de tamaño fijo, 180\,MiB. \\
Tiempo de arranque & $\leq$ 15\,s & 17\,s & Instante en que
\id{robot-server} queda activo, en la Raspberry Pi 4 con WiFi. \\
RAM en operación & --- & \pendiente{medir} & \id{/proc/meminfo} y
\id{VmRSS} de cada proceso. \\
CPU en operación & --- & \pendiente{medir} & \id{/proc/stat} durante 60\,s. \\
\bottomrule
\end{tabularx}
\end{center}

El rootfs queda en el 65\,\% del presupuesto. El arranque excede la
referencia por 2\,s: \id{robot-server} espera a que la red esté operativa
(\id{network-online.target}), y la asociación WPA2 más la concesión DHCP
dominan ese tiempo. Se considera una desviación justificada, porque sin red
el panel de control no es alcanzable y el servicio no estaría operativo en
el sentido del enunciado. El método completo está en \id{docs/metricas.md}.

\subsection{Validación del costo y del impacto}

\begin{itemize}
  \item \textbf{Costo.} \pendiente{comparar el costo real con la estimación
        de 49--99 USD.}
  \item \textbf{Impacto ambiental.} Todos los módulos quedaron con conectores
        desconectables; la Raspberry Pi se devuelve sin soldaduras ni
        adhesivos; el pack es recargable y desmontable.
\end{itemize}

\subsection{Hallazgos de la validación y acciones correctivas}
\label{sec:hallazgos}

Validar no fue solo confirmar: varias pruebas cambiaron el diseño.

\begin{longtable}{L{0.5cm}L{4.4cm}L{4.9cm}L{4.7cm}}
\toprule
& \textbf{Hallazgo} & \textbf{Causa} & \textbf{Acción} \\
\midrule
\endhead
1 & El límite de reintentos de \id{systemd} no existía. &
\id{StartLimitIntervalSec} estaba en \id{[Service]}; desde systemd 229 va en
\id{[Unit]} y en otro lugar se ignora en silencio. & Se movió a \id{[Unit]}.
Lo detectó \id{systemd-analyze verify}. \\
2 & Con los optoacopladores, los motores girarían al revés. & El PC817 en
emisor común invierte la señal. & Compensación en un único punto de
\id{lib_motors.c}; el simulador invierte igual que la placa, así que quitarla
hace fallar la prueba. \\
3 & Con PWM de 1\,kHz los motores no se movieron. & El PC817 con pull-up de
4.7\,k$\Omega$ tarda decenas de microsegundos en apagarse; sumado a los 2\,V
que cae el L298N, a ciclos medios el motor no alcanza su tensión de arranque.
& \pendiente{describir la corrección aplicada: PWM a 100\,Hz, pull-up de
1\,k$\Omega$ o velocidad fija} \\
4 & El servidor arrancaba a los 130\,s. & \id{eth0} sin cable y dos clientes
DHCP sobre \id{wlan0} agotaban la espera de red. & Ethernet marcada como no
requerida y un solo cliente DHCP: 17\,s. \\
5 & El vaivén del servo sacudía el sensor y el chasis. & El servo iba a su
velocidad máxima entre paradas. & Rampa a un tercio de la velocidad, con un
pulso intermedio cada 20\,ms. \\
6 & El \emph{build} colgaba el host al 67\,\%. & Tres compilaciones pesadas
a la vez agotaron 7\,GB de RAM y el swap. & Dos hilos y regulación por
presión de recursos (\id{BB_PRESSURE_MAX_MEMORY}). \\
\bottomrule
\end{longtable}

\subsection{Limitaciones conocidas}

\begin{itemize}
  \item El radar usa un solo sensor físico para cubrir las detecciones
        frontal y lateral. El diseño deja libres los GPIO 22 y 10 para un
        segundo HC-SR04 fijo al frente.
  \item Sin encoders, la distancia recorrida depende de la calibración del
        modelo de los motores; el mapa sirve para la sesión, no como posición
        absoluta.
  \item Una lectura sin eco no significa camino libre: una pared oblicua
        desvía el pulso. El mapa no usa esas lecturas para despejar celdas.
  \item Los requerimientos opcionales no se implementaron.
\end{itemize}

% ===========================================================================
\section{Conclusiones}

El diseño cumple los requerimientos obligatorios con un rootfs de 129\,MB y
un modelo físico construido con el inventario del grupo y una compra acotada.
Las decisiones que más pesaron ---radar sobre servo, pack 2S con dos rieles,
MPU-6050 fusionado con el modelo de los motores--- se tomaron comparando
alternativas con criterios técnicos, económicos, ambientales y sociales, y
cada una dejó una debilidad conocida que el propio diseño atiende.

La validación en tres niveles permitió avanzar el software antes de tener el
hardware, pero también mostró su límite: el problema de la PWM a través del
optoacoplador solo apareció en la placa real. El diseño de un sistema
embebido se cierra en el hardware.

\section{Referencias}

\begin{itemize}
  \item J. González Gómez, \emph{Especificación Proyecto I}, CE-1113 Sistemas
        Empotrados, Tecnológico de Costa Rica, 2026.
  \item Yocto Project, \emph{Reference Manual} y \emph{Development Tasks
        Manual}, versión 5.0 (\id{scarthgap}).
  \item Raspberry Pi Ltd., documentación de la Raspberry Pi 4 Model B: GPIO,
        \id{config.txt} y \emph{device tree overlays}.
  \item \id{pigpio}, documentación de la interfaz \id{pigpiod_if2}.
  \item Hojas de datos: Sharp PC817, STMicroelectronics L298N, InvenSense
        MPU-6050, Diodes Inc. PAM8403, Monolithic Power MP2307, HC-SR04.
  \item Repositorio del proyecto:
        \url{https://github.com/Randall-BL/EmpotradosProyecto1}.
\end{itemize}

\end{document}
